\subsection{SEMI-SIMPLE ELEMENTS AND NILPOTENT ELEMENTS IN SEMI-SIMPLE LIE ALGEBRAS}

PROPOSITION 3. Let M be a finite-dimensional vector space over K and g a semi-simple subalgebra of gl(M). Then g contains the semi-simple and nilpotent components of its elements.

If \(\mathbf{K}_1\) is an extension of \(\mathbf{K}\) , the Killing form of \(\mathfrak{g}_{(\mathbf{K}_1)}\) is the extension to \(\mathfrak{g}_{(\mathbf{K}_1)}\) of that of \(\mathfrak{g}\) ( \(\S 3\) , no. 8) and hence is non-degenerate; therefore \(\mathfrak{g}_{(\mathbf{K}_1)}\) is semi-simple. It therefore suffices to prove Proposition 3 when the base field is algebraically closed, which we shall henceforth assume to be the case.

For every subspace \(\mathbf{N}\) of \(\mathbf{M}\) , let \(\mathfrak{g}_{\mathbf{N}}\) be the subalgebra of \(\mathfrak{gl}(\mathbf{M})\) consisting of the elements which leave \(\mathbf{N}\) stable and whose restriction to \(\mathbf{N}\) has trace zero. As \(\mathfrak{g} = \mathcal{D}\mathfrak{g}\) , \(\mathfrak{g} \subset \mathfrak{g}_{\mathbf{N}}\) if \(\mathbf{N}\) is stable under \(\mathfrak{g}\) . Then let \(\mathfrak{g}^*\) be the intersection of the normalizer of \(\mathfrak{g}\) in \(\mathfrak{gl}(\mathbf{M})\) and the algebras \(\mathfrak{g}_{\mathbf{N}}\) where \(\mathbf{N}\) runs through the set of subspaces of M which are stable under g. As the semi-simple (resp. nilpotent) component s (resp. n) of \(x \in \mathfrak{gl}(M)\) is a polynomial in x with no constant term and ad s (resp. ad n) is the semi-simple (resp. nilpotent) part of ad x ( \(\S\) 5, no. 4, Lemma 2), clearly \(x \in g^{*}\) implies \(s \in g^{*}\) and \(n \in g^{*}\) ; it therefore suffices to show that \(g^{*} = g\) . Since g is a semi-simple ideal of \(g^{*}, g^{*} = a \times g\) (no. 1, Corollary 1 to Proposition 1). Let \(a \in a\) and let N be a subspace which is minimal among the non-zero subspaces of M which are stable under g. The restriction of a to N is a scalar multiple of the identity by Burnside's Theorem, has trace zero by construction and hence is zero since K is of characteristic 0. As M is the direct sum of subspaces such as N, it follows that a = 0 and hence \(g^{*} = g\) .

COROLLARY. An element \(x\) of \(\mathfrak{g}\) is a semi-simple (resp. nilpotent) endomorphism of \(\mathbf{M}\) if and only if \(\operatorname{ad}_{\mathfrak{g}} x\) is a semi-simple (resp. nilpotent) endomorphism of \(\mathfrak{g}\) .

Let \(s\) (resp. \(n\) ) be the semi-simple (resp. nilpotent) component of \(x \in \mathfrak{g}\) . Then \(s \in \mathfrak{g}\) and \(n \in \mathfrak{g}\) (Proposition 3). Then \(\mathrm{ad}_{\mathfrak{g}} s\) (resp. \(\mathrm{ad}_{\mathfrak{g}} n\) ) is the semi-simple (resp. nilpotent) component of \(\mathrm{ad}_{\mathfrak{g}} x\) , by Lemma 2 of §5, no. 4. If \(x\) is semi-simple (resp. nilpotent) so then is \(\mathrm{ad}_{\mathfrak{g}} x\) . If now \(\mathrm{ad}_{\mathfrak{g}} x\) is semi-simple (resp. nilpotent), it is equal to \(\mathrm{ad}_{\mathfrak{g}} s\) (resp. \(\mathrm{ad}_{\mathfrak{g}} n\) ) and hence \(x = s\) (resp. \(x = n\) ) since the adjoint representation of \(\mathfrak{g}\) is faithful.

DEFINITION 3. Let g be a semi-simple Lie algebra. An element x of g is called semi-simple (resp. nilpotent) if, for every g-module M of finite dimension over K, \(x_{M}\) is a semi-simple (resp. nilpotent) endomorphism of M.

PROPOSITION 4. Let \(\mathfrak{g}\) , \(\mathfrak{g}'\) be semi-simple Lie algebras and \(f\) a homomorphism of \(\mathfrak{g}\) into \(\mathfrak{g}'\) . If \(x \in \mathfrak{g}\) is semi-simple (resp. nilpotent), so is \(f(x)\) . If \(f\) is surjective, every semi-simple (resp. nilpotent) element of \(\mathfrak{g}'\) is the image under \(f\) of a semi-simple (resp. nilpotent) element of \(\mathfrak{g}\) .

If \(\rho\) is a representation of \(\mathfrak{g}'\) , \(\rho \circ f\) is a representation of \(\mathfrak{g}\) , whence the first assertion. If \(f\) is surjective, there exists a homomorphism \(g\) of \(\mathfrak{g}'\) into \(\mathfrak{g}\) such that \(f \circ g\) is the identity homomorphism of \(\mathfrak{g}'\) (no. 1, Corollary 2 to Proposition 1) and the second assertion then follows from the first.

THEOREM 3. Let g be a semi-simple Lie algebra.

\begin{enumerate}
\def\labelenumi{(\alph{enumi})}
\item
  Let \(x \in \mathfrak{g}\) . If there exists a faithful representation \(\rho\) of \(\mathfrak{g}\) such that \(\rho(x)\) is a semi-simple (resp. nilpotent) endomorphism, then \(x\) is semi-simple (resp. nilpotent).
\item
  Every element of \(\mathfrak{g}\) can be written uniquely as the sum of a semi-simple element and a nilpotent element which commute with one another.
\end{enumerate}

Suppose that the hypothesis of (a) holds. Let \(\sigma\) be a representation of \(\mathfrak{g}\) , \(\mathfrak{b}\) the supplementary ideal of the kernel of \(\sigma\) and \(\alpha\) the projection of \(\mathfrak{g}\) onto \(\mathfrak{b}\) . Then \(\mathrm{ad}_{\mathfrak{g}}x\) is semi-simple (resp. nilpotent) by the Corollary to Proposition 3 and hence \(\mathrm{ad}_{\mathfrak{b}}\alpha(x)\) is semi-simple (resp. nilpotent). As \(\sigma(x)=\sigma(\alpha(x))\) , the first assertion follows from the Corollary to Proposition 3. The second result then follows from Proposition 3 applied to a faithful representation.

